\documentclass[11pt,a4paper]{article}
\usepackage{fontspec,amsmath,amssymb,amsthm,microtype}
\setmainfont{lmroman10-regular.otf}[BoldFont=lmroman10-bold.otf,ItalicFont=lmroman10-italic.otf,BoldItalicFont=lmroman10-bolditalic.otf]
\setmonofont{lmmono10-regular.otf}[Scale=MatchLowercase]
\usepackage[margin=25mm]{geometry}
\usepackage[unicode,hidelinks]{hyperref}
\setlength{\emergencystretch}{3em}
\setlength{\parskip}{0.4em}
\setlength{\parindent}{0pt}
\allowdisplaybreaks[2]
\raggedbottom
\title{\texorpdfstring{The L\ensuremath{{}^{2}} recursion is the original implemented KLMC chain}{The L² recursion is the original implemented KLMC chain}}
\author{LangevinLean mathematical exposition}
\date{}
\begin{document}\maketitle
\noindent Source module: \texttt{Kinetic\allowbreak{}Langevin/\allowbreak{}Lp\allowbreak{}Dynamics.\allowbreak{}lean}. The accompanying HTML pairs this mathematical exposition with the exact Lean statements and records the scope of the formalization.


Analysis on equivalence classes of random vectors is useful only if it represents the intended algorithm. The lemmas here prove that the actual lifted exponential update, KLMC noise scaling and every recursively defined iterate agree almost everywhere with the original pointwise vector formulas. No Gaussian or independence assumptions are required for this algebraic correspondence.


\section{\texorpdfstring{Read L\ensuremath{{}^{2}} phase points and identify the exponential update}{Read L² phase points and identify the exponential update}}


\texttt{lp\allowbreak{}Phase\allowbreak{}Coe} reads the velocity and position representatives of a pair of L\ensuremath{{}^{2}} classes at the same sample. A private affine-operation lemma combines the almost-everywhere identities for scalar multiplication, subtraction and addition. Apply it to the velocity formula \(\psi_0v-\psi_1G(x)+\eta_v\). For position, rewrite its affine expression as \(\psi_1v-\psi_2G(x)+(x+\eta_x)\), then use the same lemma. Finally \texttt{lp\allowbreak{}Drift\_\allowbreak{}coe} identifies \(G(x)(\omega)\) with \(\nabla f(x(\omega))\). The two component equalities give equality of the actual phase update and the original pointwise \texttt{exponential\allowbreak{}Step}.



{\raggedright\noindent Lean declarations: \texttt{lp\allowbreak{}Phase\allowbreak{}Coe}, \texttt{lp\_\allowbreak{}exponential\allowbreak{}Step\_\allowbreak{}coe}.\par}

\section{\texorpdfstring{Expand affine operations on L\ensuremath{{}^{2}} representatives}{Expand affine operations on L² representatives}}


The representative of a sum or scalar multiple of Lp classes agrees with the corresponding operation almost everywhere. The private helper gathers the four rules needed for \(aX-bY+Z\) into one common full-measure set. On that set, it rewrites the representative of the outer sum, the inner difference, and the two scalar multiples. The resulting equality is precisely the pointwise affine formula. This step accounts for the fact that Lp elements are equivalence classes; pointwise equalities of their arbitrarily selected representatives cannot be assumed everywhere.



{\raggedright\noindent Lean declarations: \texttt{affine\_\allowbreak{}lp\_\allowbreak{}coe}.\par}

\section{\texorpdfstring{Retain the exact KLMC noise scaling through every iterate}{Retain the exact KLMC noise scaling through every iterate}}


The two phase components of a scaled L\ensuremath{{}^{2}} innovation agree almost everywhere with the same scalar times their representatives. Combining this with the exponential-step correspondence proves the KLMC correspondence, including its \(\sqrt{2\gamma}\) fluctuation-dissipation factor. Induction then proves the entire chain identity: at zero both chains use the supplied initial point; at a successor, intersect the almost-everywhere previous-chain identity with the current-step identity, and rewrite the recursive definitions. This shows the L\ensuremath{{}^{2}} error estimates concern the original numerical sampler for the supplied innovation stream, rather than a different abstract sequence.



{\raggedright\noindent Lean declarations: \texttt{lp\allowbreak{}Phase\allowbreak{}Coe\_\allowbreak{}smul\_\allowbreak{}ae}, \texttt{lp\_\allowbreak{}klmc\allowbreak{}Step\_\allowbreak{}coe}, \texttt{lp\_\allowbreak{}klmc\allowbreak{}Chain\_\allowbreak{}coe}.\par}

\end{document}
